\clearpage
\section{Integral Projectors, Cotangent Multiplicity, and Wiles Defect}

\subsection{Projector denominator as a congruence exponent}
Let $V=U\oplus W$ be a rational decomposition and let $\Lambda\subset V$ be an integral lattice. Let $e_U$ be the rational projector onto $U$. Define
\[
 \Lambda_U=\Lambda\cap U,
 \qquad
 \Lambda_W=\Lambda\cap W,
\]
and
\[
 C_U=\Lambda/(\Lambda_U+\Lambda_W).
\]

\begin{theorem}[Integral projector theorem]
The least positive integer $\eta_U$ satisfying
\[
 \eta_U e_U(\Lambda)\subseteq\Lambda
\]
is the exponent of the finite quotient $C_U$.
\end{theorem}

\begin{proof}
For $x=u+w\in\Lambda$, the condition $m e_U(x)=mu\in\Lambda_U$ is equivalent to
\[
 mx=mu+mw\in\Lambda_U+\Lambda_W.
\]
Thus $m$ clears every projector denominator exactly when it annihilates every class in $C_U$.
\end{proof}

For a primitive rank-one vector $v\in\Z^n$, the orthogonal projector denominator is
\[
 \langle v,v\rangle,
\]
and the congruence quotient is cyclic of that order.

\subsection{Why one cyclic quotient is not enough for split persistence}
The split persistence module
\[
 \mathcal P_c
 \cong
 \bigoplus_{p^e\parallel c}
 (\Z/p\Z)^{e-1}
\]
has large $p$-rank. The $p$-torsion of a cyclic target has rank at most one. Therefore every map from $\mathcal P_c$ to a cyclic group has a forced kernel of order at least
\[
 \prod_{p^e\parallel c}p^{\max(e-2,0)}.
\]
This is a true obstruction to the split source. It motivates the Hecke cotangent module
\[
 \mathfrak p/\mathfrak p^2,
\]
which can retain multiple independent congruence directions before cyclic compression.

\subsection{A toy Hecke order}
Consider
\[
 \mathbb T(\mathbf m)
 =
 \{(x_0,x_1,\ldots,x_r)\in\Z^{r+1}:x_i\equiv x_0\pmod{m_i}\}.
\]
For the target coordinate $x_0$, the cotangent module is
\[
 \boxed{
 \mathfrak p/\mathfrak p^2
 \cong
 \bigoplus_{i=1}^{r}\Z/m_i\Z.
 }
\]
The cyclic congruence quotient is
\[
 \boxed{
 \Z/\lcm(m_1,\ldots,m_r)\Z.
 }
\]
The quotient of cotangent order by congruence order is
\[
 \frac{\prod_i m_i}{\lcm(m_i)}.
\]
This is a toy Wiles-defect factor. It measures the multiplicity discarded by cyclic compression.

The comparison between cotangent and congruence data, including non-complete-intersection defects and local--global formulas, is developed in the Wiles-defect literature~\cite{BKM1,BKM2}.

\subsection{Complete-intersection control case}
Let $F\in\Z[X]$ be monic and squarefree over $\Q$, with an integer simple root $a_0$. Put
\[
 \mathbb T=\Z[X]/(F)
\]
and let $\lambda:\mathbb T\to\Z$ evaluate at $a_0$. The ring is finite free over $\Z$ and reduced over $\Q$. Then the congruence quotient $\Z/\lambda(\operatorname{Ann}I)$ and cotangent module both have order
\[
 |F'(a_0)|.
\]
To see this directly, write $F=(X-a_0)G$. The kernel $I$ is generated by $X-a_0$. Modulo $I^2$, the relation is $G(a_0)(X-a_0)=0$, so $I/I^2\cong\Z/F'(a_0)\Z$. The annihilator of $I$ is generated by $G$ modulo $F$, whose evaluation is $G(a_0)=F'(a_0)$. Thus the congruence ideal is $F'(a_0)\Z$ and the defect is one. Extra multiplicity is therefore tied to non-complete-intersection structure rather than being a normalization artifact.

\subsection{Two prime indices}
The persistence is indexed by conductor primes $q\mid abc$. Congruence depth is indexed by coefficient primes $\ell$. The natural local signal is a matrix
\[
 D_f(\ell,q).
\]
Let $E/\Q_q$ have split multiplicative reduction with Tate parameter $q_E$, and let $\ell\ne q$. Scalarity of the full local Galois representation on $E[\ell^n]$ forces
\[
 n\le\min\{v_\ell(q-1),v_\ell(v_q(q_E))\}.
\]
In a Tate basis, the cyclotomic graded diagonal and the inertia Kummer extension respectively force $\ell^n\mid(q-1)$ and $\ell^n\mid v_q(q_E)$~\cite{SerreAbelian1989}. This is a necessary condition, rather than a converse or a condition on Frobenius alone. These off-diagonal channels carry tame congruence information. They do not naturally reproduce the full base-$q$ factor $q^{e-1}$. This mismatch motivates the later filtered contact correction and the diagonal coefficient-prime channel.

